\section{Real-Time Guardrails}
\label{sec:realtime}

Real-time guardrails are mechanisms that must operate within the latency budget of
a live inference request.
We define this operationally: a guardrail is \emph{real-time} if it is designed to
fire before or during token delivery to the end user, rather than in an offline or
batch evaluation pipeline.
This section reviews the principal systems and identifies the open problems
that the literature has not resolved.

\subsection{Inference-Time Classification Systems}

\subsubsection{Llama Guard}

Llama~Guard~\cite{inan2023llama} is the most fully documented real-time guardrail
system in the literature.
It frames content moderation as an instruction-following classification task:
a 7B-parameter LLM fine-tuned on a safety taxonomy produces \texttt{safe} or
\texttt{unsafe} as its first output token, followed by one or more violated category
labels.
The instruction-tuning framing allows the same model to be applied to both the
user prompt (pre-generation) and the model response (post-generation) without
retraining; only the prompt template changes.

This architectural description is the only real-time guardrail claim that survived
all three self-review verification passes (of 25 extracted claims in this section)
under the protocol described in Section~\ref{subsec:verification_protocol}.
Meta has since released Llama Guard~2 and Llama Guard~3~\cite{meta2024llamaguard2},
extending the safety taxonomy to align with the MLCommons AI Safety standard.
Of direct relevance to Section~\ref{sec:multimodal}, Llama Guard~3 extends the
classification architecture to support image inputs alongside text, making it the
only documented progression of the original text-only Llama Guard design toward
multimodal safety classification.
The architectural extension is described in Meta's release materials but has not
been evaluated independently; the image-safety classification accuracy, latency
overhead of the vision pathway, and behaviour against cross-modal jailbreaks are
not established in the reproducible academic literature.
Quantitative comparisons between versions (v1, v2, v3) have not been independently
verified and are treated as unverified throughout this review.

\subsubsection{ShieldLM}

ShieldLM~\cite{zhang2024shieldlm} is a bilingual (Chinese/English) safety
detector built by instruction-tuning a 7B-parameter LLM.
The system is designed to be customisable via natural language rule specification at
inference time, enabling domain-specific policy enforcement without retraining.
Its claimed accuracy advantages over GPT-4 and Llama Guard~2, and the claimed
latency reduction from token-output truncation, did not survive adversarial
verification in this review.
The underlying architectural approach---customisable natural-language policy
specification---is a meaningful design contribution regardless of the specific
numbers, and warrants further reproducible study.

\subsubsection{WildGuard}

WildGuard~\cite{han2024wildguard} addresses three safety tasks in a single
inference pass: prompt harmfulness classification, response harmfulness classification,
and refusal behaviour detection.
This multi-task framing is architecturally significant because it avoids the latency
cost of running separate classifiers.
The claimed reduction in jailbreak attack success rate from 79.8\% to 2.4\% did not
survive adversarial verification; the comparison baseline (Llama Guard~2 leaving
53.1\% ASR) similarly could not be confirmed.
The single-pass multi-task design remains a promising direction.

\subsubsection{The Dual Reward Model Architecture}

The only quantified, reproducible arbitration mechanism in the literature is the
dual reward model design from Llama~2~\cite{touvron2023llama2}.
Two independent reward models---one for helpfulness, one for safety---are trained
separately.
At inference time, the safety reward model overrides the helpfulness reward model
when the safety score $R_s(g|p) < 0.15$, where a threshold was selected to achieve
precision 0.89 and recall 0.55 on the Meta Safety test set.
This constitutes the most concrete, verifiable description of a real-time
safety/helpfulness arbitration mechanism in the reviewed literature.

\subsection{Pipeline-Based Frameworks}

\subsubsection{NeMo Guardrails}

NeMo Guardrails~\cite{rebedea2023nemo} implements a dialogue-management-inspired
runtime that applies rails at multiple interception points in the inference pipeline:
input, knowledge retrieval, dialogue management, action execution, and output.
Rails are specified in Colang, a domain-specific language designed to express
dialogue flows and policy constraints.
The framework has undergone substantial development since the founding 2023 preprint,
with multiple major releases through 2024--2025 adding support for additional LLM
backends, expanded rail types, and updated streaming semantics; readers consulting
the codebase or documentation should be aware that the implementation may differ
materially from what the 2023 paper describes.
In accordance with the review's scope and temporal specificity standard, all
analytical statements below pertain to the documented 2023 description only.

The architectural motivation---decoupling safety policy from model weights so that
rails can be updated without retraining---is sound and represents an important
engineering contribution.
However, specific claims about pipeline stages, the number of LLM round-trips
per user turn, and the timeline of streaming support did not survive author-verified
scrutiny.
A key limitation noted in the caveats of the original paper but underemphasised
in subsequent citations is that each rail invocation may trigger additional LLM
calls, adding cumulative latency that is difficult to bound in the general case.

\subsubsection{Guardrails AI}

Guardrails AI (guardrailsai.com) provides an open-source framework for defining
input and output validators (``guards'') around LLM calls.
Guards can invoke classifiers, regex patterns, or LLM-based validators.
The framework has no peer-reviewed evaluation of its safety guarantees or latency
characteristics in the academic literature.
From the standpoint of this review, the analytically relevant point is the absence
of published empirical evidence: whether Guardrails AI performs well in production
is unknown because the data required to assess that question have not been made
available in reproducible form.

\subsection{The Streaming Safety Gap}
\label{subsec:streaming_gap}

The most significant unresolved problem in real-time guardrails is \emph{streaming
safety}: how to enforce safety constraints when tokens are delivered to the user
incrementally, before a complete response is available for inspection.

Post-generation classifiers are structurally incompatible with streaming delivery:
by the time the classifier fires, the harmful content has already reached the user.
Pre-generation prompt classifiers address only the input side and miss
response-level harms generated from benign inputs.
In-generation safety checking---monitoring the growing token prefix as generation
proceeds---remains an open engineering problem with no verified solution in the
literature.

\paragraph{Relation to constrained-generation literature}
The adjacent field of constrained text generation has produced prefix-level
decoding techniques that enforce hard constraints during generation.
FUDGE~\cite{yang2021fudge} trains a future discriminator that re-weights the
next-token distribution to steer generation toward a target attribute; COLD
Decoding~\cite{qin2022cold} uses Langevin dynamics to satisfy multiple
differentiable constraints jointly; and NeuroLogic
decoding~\cite{lu2021neurologic} enforces predicate-logic constraints over
token sequences via beam search with constraint tracking.
Speculative decoding~\cite{leviathan2023fast} offers a complementary
architectural pattern in which a small draft model generates candidate tokens
that a larger model verifies, creating a natural slot for a concurrent safety
check in the verification step.

These techniques are \emph{adjacent} to, but do not resolve, the streaming
safety problem for three reasons.
First, constrained-generation methods assume that the constraint function is
differentiable or efficiently computable over short prefixes; safety is
a semantic property that may require inspecting long context windows to evaluate
reliably, making prefix-level approximations prone to high false-negative rates
on adversarially crafted prompts.
Second, FUDGE, COLD Decoding, and NeuroLogic all introduce per-token or per-step
overhead (additional forward passes or energy-function evaluations) that compounds
at the token level and may push end-to-end latency beyond acceptable bounds for
production streaming deployments.
Third, none of these methods has been evaluated against adaptive adversaries who
know the constraint mechanism and can craft inputs to minimise the safety
signal on partial prefixes---the precise threat model of interest here.
In-generation safety checking therefore remains open, and the constrained-generation
literature, while offering useful structural inspiration, does not transfer directly
to this setting without addressing these three gaps.

\paragraph{Null finding (documented corpus entry)}
A systematic search of the corpus for papers reporting empirical prefix-level safety
classification — that is, studies that train or evaluate a classifier on token
prefixes of LLM outputs for the purpose of safety detection — returned zero results.
No paper in the 75-paper corpus, and no paper identified but excluded during
full-text review, has empirically characterised the relationship between prefix
length and safety detection accuracy for any harm category.
This null finding is itself a substantive result: the streaming safety problem is
not merely unsolved but unstudied at the empirical level, leaving open questions
RT1--RT3 (Section~\ref{sec:open_problems}) entirely without quantitative bounds.
The closest adjacent evidence comes from the constrained-generation literature
reviewed above: FUDGE~\cite{yang2021fudge} operates on prefixes of minimum length
determined by the future discriminator's receptive field (roughly 10--20 tokens
in the reported experiments), suggesting that very short prefixes may be
insufficient for semantic safety decisions; however, FUDGE's discriminator targets
stylistic attributes, not adversarial harmful intent, and the analogy is
therefore only weakly informative.

Key open questions within this gap include:

\begin{itemize}
  \item What is the minimum prefix length at which a safety classifier can reliably
        detect harmful generation trajectories, and how does this vary by harm category
        and model?
        (No empirical lower bound currently exists in the literature.)
  \item Can a safety classifier be parallelised with the generation process without
        introducing head-of-line blocking?
        (Speculative decoding~\cite{leviathan2023fast} provides a structural template
        but has not been adapted for safety classification.)
  \item What is the cost (in terms of both compute and over-refusal rate) of
        conservative early-stopping on ambiguous prefixes?
        (No study has measured the utility penalty of prefix-level early stopping
        under realistic streaming latency constraints.)
\end{itemize}

\subsection{Evaluation Against Adaptive Adversaries}

A significant gap in the deployment-relevance axis of the quality assessment
(Table~\ref{tab:quality}) is that none of the five real-time systems reviewed above
has been evaluated against an adaptive adversary who knows the guardrail mechanism
and optimises inputs to evade it.
The token-level adversarial suffix literature (GCG attack,~\cite{zou2023universal})
demonstrates that gradient-based optimisation can reliably find suffixes that cause
aligned LLMs to produce harmful outputs; these attacks are directly relevant to
classifier-based guardrails because a suffix that causes the LLM to begin generating
harmful content may also suppress the safety signal in a prefix-monitoring classifier
that reads from the same token stream.
Circuit Breakers~\cite{zou2024improving} represent a complementary weight-level
defence that attempts to disrupt the internal representations that harmful generation
relies on, rather than filtering outputs post-hoc; its evaluation includes white-box
adaptive attacks, making it the strongest adaptive-adversary robustness result in the
corpus, though its quantitative claims did not survive author-verified scrutiny in
this review.
No paper in the corpus reports evaluating the inference-time classifier systems
(Llama Guard, ShieldLM, WildGuard, or NeMo Guardrails) against GCG-style adaptive
suffix attacks.
This omission is methodologically significant: a guardrail evaluated only against
non-adaptive attacks provides no assurance of robustness once the attack distribution
shifts to include adversaries who optimise against the specific classifier.
The absence of adaptive evaluation is therefore noted as a cross-cutting
deployment-relevance gap affecting all five systems reviewed in this section.

\subsection{Latency-Accuracy Tradeoffs}

Table~\ref{tab:latency} summarises the latency and accuracy figures \emph{as
reported} by each system's authors.
All figures are marked as unverified (U): they are reproduced from the source
papers but could not be independently confirmed under the adversarial verification
protocol (Section~\ref{subsec:verification_protocol}).
They are included to give practitioners a reference point while making the
reproducibility status explicit.

\begin{table}[t]
  \caption{Reported latency and accuracy figures for real-time guardrail systems.
           All entries are \textbf{unverified (U)}: figures are taken from author
           self-reports and could not be independently confirmed.
           N/A = not reported.
           \textbf{The latency column is not comparable across rows}: hardware,
           batch size, concurrency, and network conditions differ for every system.
           A hardware column is included to make this explicit; do not rank systems
           by the latency column alone.}
  \label{tab:latency}
  \centering
  \small
  \begin{tabular}{p{2.5cm}p{1.8cm}p{2.4cm}p{3.4cm}p{1.5cm}p{0.8cm}}
    \toprule
    \textbf{System} & \textbf{Hardware} & \textbf{Latency (rep.)} & \textbf{Accuracy (reported)} & \textbf{Stream?} & \textbf{St.} \\
    \midrule
    \multicolumn{6}{l}{\textit{Hardware-specified (non-comparable with each other or with rows below):}} \\
    Llama Guard~\cite{inan2023llama}       & A100        & ${\sim}200$ ms   & F1 0.82 (OpenAI Mod.)    & No      & U \\
    ShieldLM~\cite{zhang2024shieldlm}      & V100        & ${\sim}150$ ms   & Acc. 90.6\% (bilingual)  & No      & U \\
    \midrule
    \multicolumn{6}{l}{\textit{Hardware unspecified:}} \\
    WildGuard~\cite{han2024wildguard}      & ---         & N/A              & ASR 2.4\% (claimed)      & No      & U \\
    NeMo Guardrails~\cite{rebedea2023nemo} & ---         & 1--3$\times$ LLM & N/A                      & Partial & U \\
    Aegis~\cite{su2024aegis}               & ---         & N/A              & F1 0.84 (ensemble)       & No      & U \\
    Azure Content Safety~\cite{azure2024contentsafety} & Cloud API & N/A   & N/A (proprietary)        & No      & U \\
    \bottomrule
  \end{tabular}
  \par\smallskip
  \noindent\small Subgroups make explicit that the A100 and V100 figures
  (Llama Guard, ShieldLM) are on incompatible hardware baselines, and that all
  remaining rows have no stated hardware; no latency comparison across any rows
  is meaningful.
\end{table}

Three architectural families represent distinct design intents in this space.
Because all entries in Table~\ref{tab:latency} are unverified and were obtained
on incomparable hardware, the following characterisations describe the
\emph{design rationale} of each family rather than making comparative performance
claims:

\begin{description}
  \item[LLM-as-judge] A second, full-size LLM evaluates each response.
        The design intent is maximum accuracy by leveraging full model reasoning,
        at the cost of high latency; this architecture is structurally incompatible
        with streaming delivery regardless of hardware, because a full response must
        be available before the judge can be invoked.
  \item[Classifier-based] A smaller, purpose-trained classifier (BERT-scale to 7B)
        evaluates each response.
        The design intent is reduced latency at the cost of reduced accuracy;
        this family is \emph{potentially} adaptable to streaming with prefix-based
        invocation, but no verified implementation has demonstrated this in practice.
  \item[Hybrid] A lightweight classifier gates whether a full LLM judge is invoked.
        The design intent is to achieve a better efficiency frontier by reserving
        the expensive judge for ambiguous cases; however, this introduces
        architectural complexity and dual failure modes (both the classifier and the
        judge must be correct for end-to-end safety).
\end{description}

No verified quantitative evidence exists in the literature that any of these
families dominates the others on any reproducible latency-accuracy metric.
The latency-accuracy design space for runtime guardrail architectures therefore
remains empirically uncharted in reproducible literature, representing a direct
opportunity for future work.

\subsection{Formal Guarantees}

The Aegis system~\cite{su2024aegis} proposes an online adaptive ensemble of safety
expert models with a theoretical no-regret bound.
The specific mathematical claims (regret $O(\sqrt{T \log K})$, phased
adaptation/compliance rounds) did not survive adversarial verification for two
reasons: (i) the proof of the regret bound references an appendix that was not
present in any accessible version of the preprint at the time of review; and
(ii) the definition of a ``compliance round'' in the phased-adaptation scheme is
underspecified, making it impossible to verify that the empirical algorithm
described in the paper matches the theoretical model for which the bound is derived.
Beyond these proof-level defects, it is worth noting---as developed in
open problem FG1 (Section~\ref{sec:open_problems})---that the regret formulation
itself may be inappropriate for the safety setting regardless of proof correctness:
the regret bound assumes a stationary sequence of adversarial inputs against which
the ensemble adapts, but a real jailbreak adversary is non-stationary and will
adapt to the deployed classifier, invalidating the bound after the first adaptive
attack iteration.

Kumar et al.~\cite{kumar2023certifying} attempt to provide formal safety
guarantees against adversarial prompting via an erase-and-check procedure that
provides probabilistic certificates on whether a prompt can be detected as
adversarial within a given perturbation radius.
This claim did not survive adversarial verification in this review for a specific
reason: the certification guarantee is conditioned on the adversary's perturbation
being bounded by a fixed token-deletion radius, a constraint that is not satisfied
by the free-form jailbreaks (multi-turn priming, typographic embedding, cipher
substitution) that dominate the empirical attack literature.
The formal guarantee is therefore valid within a narrow, largely hypothetical
threat model and does not generalise to the realistic deployment threat model
characterised in Section~\ref{sec:background}.

No verified formal guarantees exist in the current literature for:
(i) the false-negative rate of any runtime guardrail against novel jailbreaks
not present in the training distribution; or
(ii) the soundness of any streaming safety mechanism under adaptive adversaries.
This is a foundational theoretical gap.
